\documentclass[a4paper]{article}     
\usepackage{amsfonts}
\usepackage{amsmath}
\usepackage{amssymb}
\usepackage{graphicx}

\begin{document}

\noindent \large Auteur: Xander De Blauwer \\
\noindent \large Studierichting: Informatica\\
\noindent \large Oefeningenreeks 2E nr. 26.3\\

\medskip

\normalsize


$\log_x 3 = \log_3 x$ \\

$\Leftrightarrow \log_x 3 = \dfrac{\log_3 3}{\log_3 x}$ \\

$\Leftrightarrow \dfrac{\log_3 3}{\log_3 x} = \log_3 x$ (want $\log_x 3 = \log_3 x)$\\

$\Leftrightarrow \log_3^2 x = 1$ \\

Er zijn 2 mogelijke oplossingen: \\

$\log_3 x = 1 \Leftrightarrow x = 3 $\\

$\log_3 x = -1 \Leftrightarrow x = \dfrac{1}{3}$ \\






\begin{thebibliography}{999}
%\bibitem{a} Referentie
\end{thebibliography}
\end{document} 
